\vspace{-0.4cm}
\section{Methodology}
\vspace{-0.1cm}

\begin{table}
\centering
\footnotesize
\caption{Hardware configurations}
\label{tab:hw-config}
\begin{tabular}{c;{1pt/1pt}c|c;{1pt/1pt}c} 
\hline
Systolic array size   & 16x16 PEs & Memory bandwidth        & 936GB/s   \\ 
\hdashline[1pt/1pt]
\# of systolic arrays & 1024      & On-chip memory & 28MB      \\ 
\hdashline[1pt/1pt]
Frequency             & 500 MHz   & Peak perf.(MX9)   & 262 TOPS  \\
\hline
\end{tabular}
\end{table}

% % table for HW configuration
% \begin{table}[t] 
% \centering
% \footnotesize
% \caption{Hardware configurations}
% \label{tab:hw-config}
% \begin{tabular}{c|c} 
% \hline
% Systolic array size   & 16x16 PEs   \\ 
% \hdashline[1pt/1pt]
% \# of systolic arrays    & 1024         \\ 
% \hdashline[1pt/1pt]
% Frequency             & 500MHz     \\ 
% \hdashline[1pt/1pt]
% Memory bandwidth      & 936GB/s     \\ 
% \hdashline[1pt/1pt]
% On-chip memory capacity & 28MB        \\ 
% \hdashline[1pt/1pt]
% Peak performance(MX9) & 262 TOPS  \\
% \hline
% \end{tabular}
% \end{table}

\niparagraph{Models and datasets.}
For evaluation, we use DiT-XL~\cite{dit}, Pixart-$\Sigma$~\cite{pixart-sigma}, and Stable Diffusion 3 (SD3)~\cite{sd3}.
%
We denote the generated image resolution by appending $\{256, 512, 1024\}$ to the model name.
%
We use the default settings for the guidance scale of 4.0 and 5.0 for DiT-XL and Stable Diffusion 3, respectively.
%
We perform inference for 25 timesteps across all models.
%
To evaluate the quality, we employ ImageNet-val-5k for DiT-XL, and MS COCO-1k for Pixart-$\Sigma$ and SD3. 
%
We adopt three image quality metrics: fidelity with Frechet Inception Distance (FID), diversity with Inception Score (IS), and prompt alignment with CLIP Score. 

\niparagraph{Baselines.}
%
% We employ Q-DiT~\cite{qdit} and ViDiT-Q~\cite{viditq} on NVIDIA RTX-3090 as baselines.
We employ Q-DiT~\cite{qdit} and ViDiT-Q~\cite{viditq} as quality baselines, both of which are prior works using INT quantization for diffusion transformers.
%
We set the precision to 6 bits for both weights and activations to demonstrate their limitations in low-precision activation quantization.
%
As speedup baselines, we use the original FP16 model and the ViDiT-Q W8A8 model on RTX 3090.
% 
Since ViDiT-Q provides an INT8 model only for Pixart-$\Sigma$, we conduct our comparison with ViDiT-Q exclusively on Pixart-$\Sigma$.
% 

\niparagraph{Implementation.}
%
We implement the MX quantized model with PyTorch 2.0, Huggingface Diffusers library, and triton language~\cite{triton-language}. 
%
We develop our accelerator on top of the open-source DaCapo accelerator simulator~\cite{dacapo}.
%
% We set our simulator setting to match the MX9 peak performance to the GPU FP16 performance.
%
Table~\ref{tab:hw-config} shows the hardware configuration for the simulator.
We set the group size to 16 and the subgroup size to 2 for MX, and employ a batch size of 1 throughout the experiments.


